\begin{frame}{세 가지를 한꺼번에 보는 손계산}
  \begin{enumerate}
    \item 5$\times$5 입력이 3$\times$3으로 ``줄어든'' 것은
      $n_{\text{out}} = n-k+1$의 \textbf{자연스러운 결과} (다음 절에서 정리)
    \item \textbf{어느 위치에서든} 같은 가중치가 같은 패턴에 반응
      --- 파라미터 공유의 구체적 모습
    \item ``수직선''이 \textbf{원본 스케일 그대로} 다음 층으로 전달 ---
      원본이 5칸짜리 선이든 500칸짜리 선이든, 출력 지도 위에서는
      ``3$\times$3 지도의 중심 열이 밝다''로 \textbf{똑같이} 표현됨
  \end{enumerate}
  \vskip0.6em
  \begin{block}{``규모를 초월한'' 특징의 기반}
    마지막 점이 합성곱이 스케일에 무관한 특징을 만드는 기반이며,
    10.2$\sim$10.3에서 풀링·수용장이 이 아이디어를 확장한다.
  \end{block}
\end{frame}
